Differential testing of a code change is informative when the old and new versions disagree, but an empty difference set is often read as reassurance although it only reports what the chosen inputs showed. We study the reliability of such no-divergence verdicts. We define their false-assurance risk, $\mathcal{R}=P(F{=}1\mid\text{no divergence observed})$, as a selective risk relative to a declared observation boundary and adjudication budget, not an equivalence claim. SBG~V8 estimates this risk from a change-conditioned evidence profile (activation, coverage, multiplicity, reproducibility, provenance, integrity), abstains on inadequate evidence, and collects evidence through an authenticated, sealed pipeline. An Evidence Response Signature specifies eight properties a risk estimator should satisfy, tested by nine controlled evidence perturbations. On 312 executed development pairs (258 with construction labels), plain differential testing had an observed false-assurance risk of 0.169, and the predecessor's sufficiency gate 0.213 at a third of the coverage. As a label-free conformance test, the signature exposed a profile defect (irrelevant inputs lowered estimated risk in 109 of 227 cases; since fixed) and showed that the predecessor's verdicts ignore integrity failures. As predictive features it was not shown to help: at matched coverage, full calibration differed from coverage-only calibration by 1.6 percentage points (95\% interval $-6.6$ to $+16.3$), with unstable sign across seeds; only differences of about 16 points were detectable. Adversarial testing showed fail-safe abstention under result-channel and keyless post-hoc tampering, detection-only coverage of in-process forgery and evidence substitution, and one open sensing channel. Independent adjudication remains pending.
